\documentclass[10pt,a4paper]{amsart}
\usepackage[margin=19mm]{geometry}
\usepackage{amsmath,amssymb,mathtools}
\usepackage[scr=rsfs,cal=euler]{mathalfa}
\usepackage{xfrac}
\usepackage[unicode,hidelinks,bookmarks=false]{hyperref}
\newcommand{\ie}{i.e.\ }
\newcommand{\cf}{cf.\ }
\newcommand{\resp}{respectively\ }
\newcommand{\Subsection}{\subsection}
\providecommand{\nobreakdash}{\nobreak\hbox{-}}
\setlength{\emergencystretch}{2em}
\multlinegap0pt
\usepackage{bm}
\usepackage{bbm}
\usepackage{dsfont}
\usepackage{xspace}
\usepackage{cancel}
\usepackage{mathtools}
\usepackage{amsthm}
\makeatletter
\@ifundefined{prop}{\newtheorem{prop}{Proposition}}{}
\makeatother
\newcommand{\gM}{{\mathfrak M}}
\title[Congruences for traces of singular moduli and Hurwitz - Kron]{Congruences for traces of singular moduli and Hurwitz - Kronecker class numbers}
\author{Pavel Guerzhoy}
\date{Source: arXiv:2602.19544v1 (CC BY 4.0)}
\begin{document}
\maketitle

\noindent\emph{Source and licence.} Congruences for traces of singular moduli and Hurwitz - Kronecker class numbers. Version arXiv:2602.19544v1 (CC BY 4.0), distributed under the Creative Commons Attribution 4.0 International licence, as recorded in the arXiv licence metadata for this exact version and shown on the abstract page https://arxiv.org/abs/2602.19544v1. Full rights notice: licenses/arxiv-2602.19544-CC-BY-4.0.txt.\par\smallskip

\noindent\textit{Editorial scope.} Counted unit: the Prop whose source label is \texttt{prop\_lim\_jenkins} in the article (source0.tex lines 607--633, source label \texttt{prop\_lim\_jenkins}), delivered together with the article's own complete original proof of it (source0.tex lines 635--673, one \texttt{proof} environment of 39 lines). No mathematics is written, repaired or re-derived: statement and proof are the article's own text line for line. Required context retained uncounted: eq\_Kohnen\_plus (lines 236--238); eq\_jenkins (lines 598--602). Every definition, display and label of those lines is retained unchanged. Omitted from this package and not redistributed: 1--235, 239--597, 603--606, 634--634, 674--895. Nothing inside a retained excerpt is omitted or altered: every retained body is delivered exactly as the article's lines are written, so no omission receipt of any kind accompanies this package. Citations: the retained excerpts carry no citation command at all, so no bibliography entry of the article is transcribed and no reference is redistributed. Reference adaptation: none was needed and none was made; every reference inside the delivered unit resolves inside it, so no unresolved reference or citation is produced. Printed slips kept literal: no printed word, symbol, hypothesis or number of the counted statement or of its proof was altered. Layout and class: the article's own class is \texttt{amsart} with packages including amsmath, amssymb, xy, fancyvrb, hyperref, stmaryrd, amsthm, newtxtext, newtxmath; the wrapper is the lane's pinned \texttt{amsart} editorial wrapper at 10pt on A4, which restates the theorem-like environments the retained text needs and adds the article's own macro definitions that the retained text needs (\texttt{\textbackslash gM}), with extra wrapper packages (bm, bbm, dsfont, xspace, cancel, mathtools). The delivered numbering is the unit's own. Assets: no retained span contains a figure environment, a graphics-inclusion command, a PDF-page inclusion, a table or a TikZ picture; no image file is copied into this package, the package declares no asset dependency and no image file is redistributed with it. Licence: version arXiv:2602.19544v1 of this article is distributed under the Creative Commons Attribution 4.0 International licence, as recorded in the arXiv licence metadata for this exact version and shown on the abstract page https://arxiv.org/abs/2602.19544v1. A full rights notice is delivered at licenses/arxiv-2602.19544-CC-BY-4.0.txt. Characters: every retained excerpt is pure ASCII, so the delivered engine, XeLaTeX with its default Latin Modern text font, reports no missing character.
\begin{equation} \label{eq_Kohnen_plus}
n \equiv 1,2 \pmod{4} ~~ \text{implies} ~~ a(n) = 0.
\end{equation}

\begin{multline} \label{eq_jenkins}
B(D, p^{2n}d) = p^nB(p^{2n} D,d) \\
+ \sum_{k=0}^{n-1} \left( \frac{D}{p} \right)^{n-k-1} \left( B(D/p^2, p{^{2k}}d) - p^{k+1}B(p^{2k}D  , d/p^2)  \right) \\ 
+  \sum_{k=0}^{n-1} \left( \frac{D}{p} \right)^{n-k-1}  \left( \left( \left( \frac{D}{p} \right) - \left( \frac{-d}{p} \right) \right) p^k B(p^kD,d) \right)
\end{multline}

\begin{prop} \label{prop_lim_jenkins} 


Let $D=p^{2t}D_0$ with $p^2 \nmid D_0$, and
let $\epsilon = - \left( \frac{D_0}{p} \right)$.  

\begin{itemize}

\item[\normalfont(i)]
The $p$-adic limit 
\[
\tilde{g}_D = \lim_{n \rightarrow \infty} (-1)^{n\epsilon(\epsilon +1)/2}g_D|U^{2n} 
\]
exists.

\item[\normalfont(ii)]
If $\epsilon \neq 0$, then $\tilde{g}_D  \in \gM^\epsilon$.

\item[\normalfont(iii)]
If $p|D_0$ then $\tilde{g}_D = 0$.

\item[\normalfont(iv)]
If $p^2|D$,  then we have the reduction $g_D|U^{2n} \equiv g_{D/p^2}|U^{2n-2} \pmod {p^{n-1}}$.

\end{itemize}

\end{prop}

\begin{proof}

We start with the observation that, by definition, $m = 1,2 \pmod{4}$ implies $B(D,m) =0$, and, since $p^2 \equiv 1 \pmod 4$, 
the limit series $\tilde{g}_D$, if exists, satisies \eqref{eq_Kohnen_plus}.

In order to prove (iv), assume that $p^2|D$ and observe that   \eqref{eq_jenkins} implies, since $\left(\frac{D}{p}\right)=0$, that
\[
B(D,p^{2n}d) \equiv B(D/p^2,p^{2n-2}d) - \left( \frac{-d}{p} \right) p^{n-1} B(p^{2k}D,d) \pmod{p^n}.
\]
Statement (iv) follows from that, and induction in $t$ now allows us to prove statements (i), (ii), and (iii) under the assumption that $t=0$ (i.e. $p^2 \nmid D$).

%We analyze the possible cases separately.
We now assume that $p||D$. Thus $\epsilon =-\left( \frac{D}{p} \right) =0$, and the sums in  \eqref{eq_jenkins}  degenerate to single terms implying that
\[
B(D,p^{2n}d) \equiv - \left( \frac{-d}{p} \right)p^{n-1} B(p^{2k}D,d) \pmod{p^n}.
\]
Thus $\lim_{n \rightarrow \infty} B(D,p^{2n}d) =0$ for every $d$.
Assertion (i) in the case $\epsilon =0$ and assertion (iii) follow from that.

We now consider the two cases $\left(\frac{D}{p} \right) =\pm 1$ separately. 

Assume firstly that $\left(\frac{D}{p} \right) =1$ (thus $\epsilon=-1$). Then \eqref{eq_jenkins} becomes
\begin{multline} \label{eq_B_1_lim}
B(D, p^{2n}d) = p^nB(p^{2n}D,d) \\ - \sum_{k=0}^{n-1} p^{k+1}B(p^{2k}D,d/p^2) + \sum_{k=0}^{n-1} \left(1- \left( \frac{-d}{p} \right) \right) p^k B(p^{2k}D,d),
\end{multline}
and implies that $\lim_{n \rightarrow \infty} B(D,p^{2n}d)$ exists for every $d$, therefore the limit $\tilde{g}_D= \lim_{n \rightarrow \infty} g_D|U^{2n}$ exists.
Furthermore, if $ \left( \frac{-d}{p} \right) = 1$ then \eqref{eq_B_1_lim} becomes $B(D, p^{2n}d) = p^nB(p^{2n}D,d)$, therefore $\lim_{n \rightarrow \infty} B(D, p^{2n}d) =0$,
in other words, $\tilde{g}_D \in \gM^-$ as required. That finishes the proof of (i) and (ii) under the assumption that $\left(\frac{D}{p} \right) =1$.


Assume now that  $\left(\frac{D}{p} \right) =-1$ (thus $\epsilon=1$).  Then \eqref{eq_jenkins} becomes
\begin{multline*} 
B(D, p^{2n}d) =  p^nB(p^{2n}D,d)   \\ - \sum_{k=0}^{n-1} (-1)^{n-k-1}p^{k+1}B(p^{2k}D,d/p^2)  -  \sum_{k=0}^{n-1}  (-1)^{n-k-1} \left( 1 + \left( \frac{-d}{p} \right) \right) p^k B(p^{2k}D,d) \\
\equiv (-1)^{n} \sum_{k=0}^{n-1} (-1)^k p^k \left( pB(p^{2k}D,d/p^2) + \left( 1+ \left(\frac{-d}{p} \right) \right)B(p^{2k}D,d) \right) \pmod{p^n},
\end{multline*}
and implies that $\lim_{n \rightarrow \infty} (-1)^nB(D,p^{2n}d)$ exists for every $d$, therefore $\tilde{g}_D= \lim_{n \rightarrow \infty} (-1)^ng_D|U^{2n}$ exists.
Furthermore,  if $ \left( \frac{-d}{p} \right) = -1$, then again $B(D, p^{2n}d) = p^nB(p^{2n}D,d)$, therefore $\lim_{n \rightarrow \infty} B(D, p^{2n}d) =0$,
in other words, $\tilde{g}_D \in \gM^+$ as required. That finishes the proof of (i) and (ii) under the assumption that  $\left(\frac{D}{p} \right) =-1$.
\end{proof}

\end{document}
